% =========================
% puzzles/pXX.tex
% Final - The Brass Protractor Rite
% =========================

\PuzzleChapterIntro{The Brass Protractor Rite}{Classical Construction \textbullet\ Tangents \textbullet\ Reflections}{This is a rigorous construction puzzle. You must follow the straightedge-and-compass instructions precisely to generate a specific geometric figure, then calculate the sum of five resulting angles.}


\Lore{

The archivists recovered no diagram—only a sequence of geometric rites etched into vellum and sealed in brass.

Each segment, arc, and reflection had been chosen not for ornament, but for invocation. The angles were not decoration. They were a signature.

They say Athena devised the ritual herself—not to guard the Vault, but to remind it what precision feels like.

To unlock the seal, one must reconstruct her rite and calculate the sacred total.

}

\Problem

All angles \(\angle XYZ\) denote the \emph{smaller} angle at vertex \(Y\) (between \(0^\circ\) and \(180^\circ\)).

A courtyard designer performs the following straightedge-and-compass construction:

\begin{enumerate}
  \item Begin with a square labeled clockwise as \(S_1, T_1, U_1, R_1\).
  \item On side \(R_1U_1\), construct an equilateral triangle \(\triangle R_1V_1U_1\) with apex \(V_1\) chosen \emph{outside} the square.
  \item Draw the circle centered at \(R_1\) with radius \(R_1U_1\).
  \item Using segment \(R_1V_1\) as a side, construct a regular pentagon
        \(R_1V_1W_1Z_1A_2\) in clockwise order, chosen so that the pentagon lies in the half-plane
        determined by line \(R_1V_1\) \emph{opposite} \(U_1\).
  \item Since \(A_2\) lies on the circle centered at \(R_1\), draw the tangent to that circle at \(A_2\).
        Let this tangent meet the extension of line \(Z_1W_1\) at \(B_2\).
  \item Draw the circle centered at \(S_1\) with radius \(S_1T_1\).
  \item Let \(C_2\) be the intersection point of the two circles (centered at \(R_1\) and \(S_1\), equal radii)
        that lies \emph{outside} the square and in the same open half-plane bounded by line \(R_1S_1\) as \(Z_1\).
        Draw the tangent to the circle centered at \(S_1\) at \(C_2\).
        Let this tangent meet line \(R_1A_2\) at \(D_2\).
  \item Let the same tangent line from \(C_2\) meet line \(V_1W_1\) at \(E_2\).
  \item Let \(M_2\) be the intersection point of the two tangents (the tangent at \(A_2\) to the circle centered at \(R_1\),
        and the tangent at \(C_2\) to the circle centered at \(S_1\)).
  \item Reflect point \(B_2\) across the tangent line at \(C_2\) (i.e., across the line \(C_2D_2\)), and call the image \(G_2\).
  \item Reflect point \(A_2\) about point \(C_2\) (a half-turn about \(C_2\)), and call the image \(H_2\).
\end{enumerate}

\VaultDirective{Compute
\[
\angle Z_1B_2A_2\;+\;\angle D_2M_2A_2\;+\;\angle D_2E_2V_1\;+\;\angle R_1G_2U_1\;+\;\angle S_1C_2H_2,
\]
and report the value to the \emph{nearest whole degree}.}

\SolutionPage
\begingroup\small
Fix the square side to \(1\) with \(S_1(0,0)\), \(T_1(1,0)\), \(U_1(1,1)\), \(R_1(0,1)\). The equilateral apex is \(V_1(\tfrac12,\ 1+\tfrac{\sqrt3}{2})\). The circles meet at \(C_2(-\tfrac{\sqrt3}{2},\ \tfrac12)\); the tangent there is \(-\sqrt3x+y=2\).

Constructing the regular pentagon \(R_1V_1W_1Z_1A_2\) determines the tangent at \(A_2\) and the lines \(Z_1W_1\) and \(V_1W_1\). Solving for the intersections yields \(B_2, D_2, E_2, M_2\). Reflections yield \(G_2\) and \(H_2\). Calculating vector angles gives:
\[
\angle Z_1B_2A_2=54^\circ,\quad
\angle D_2M_2A_2=18^\circ,\quad
\angle D_2E_2V_1=72^\circ,\quad
\angle S_1C_2H_2=51^\circ,
\]
and the irrational angle \(\angle R_1G_2U_1\approx 81.9786^\circ\).
The sum is \(54+18+72+51+81.9786\dots = 276.9786\dots^\circ\), which rounds to \textbf{277}.
\endgroup

\begin{center}
\includegraphics[width=0.75\linewidth]{figures/p09_277_geometry.png}
\end{center}

\FinalAnswer{277}

\NextPage